% 本檔由 tools/build_sources.py 自各章產生，每次建置整份重寫，不要手改。
\chapter{各節出處一覽}
\label{app:sources}

表~\ref{tab:sources}依各部之序，列出各節所據之出處：第二至第四部各節先列該部所讀之本文及該節所述之刊本頁次，其次為該節另引之文獻，依第一次引用之先後排列，頁碼照該節所引。完整書目見參考文獻。

\begin{small}\setstretch{1.3}
\setlength{\LTpre}{6pt}
\def\LTcaptype{table}
\begin{longtable}{@{}>{\raggedright\arraybackslash}p{10mm}>{\raggedright\arraybackslash}p{62mm}>{\raggedright\arraybackslash}p{\dimexpr\linewidth-72mm-4\tabcolsep\relax}@{}}
\caption{各節之出處}\label{tab:sources}\\
\toprule
節 & 標題 & 出處 \\
\midrule
\endfirsthead
\toprule
節 & 標題 & 出處 \\
\midrule
\endhead
\bottomrule
\endlastfoot
\multicolumn{3}{@{}l}{導論：三篇的出處與問題} \\* \addlinespace[2pt]
\ref{sec:intro-1} & 三篇的出處與範圍 & --- \\
\ref{sec:intro-2} & 三篇各自處理的問題 & \citeauthor{alexander2000}，\citetitle{alexander2000} \newline \citeauthor{alexander1997}，\citetitle{alexander1997} \newline \citeauthor{kramer2012}，\citetitle{kramer2012} \newline \citeauthor{posner2012}，\citetitle{posner2012} \newline \citeauthor{thayer1893}，\citetitle{thayer1893} \newline \citeauthor{ferejohn2002}，\citetitle{ferejohn2002} \\
\ref{sec:intro-3} & 與 W08 的銜接 & \citeauthor{tushnet2003}，\citetitle{tushnet2003}，註 209 \newline \citeauthor{tushnet1998}，\citetitle{tushnet1998} \newline \citeauthor{tushnet1999}，\citetitle{tushnet1999} \\
\addlinespace[6pt]
\multicolumn{3}{@{}l}{Alexander 與 Schauer〈Defending Judicial Supremacy: A Reply〉} \\* \addlinespace[2pt]
\ref{sec:AS-1} & 導言：先前論點、批評與本文目的 & \citeauthor{alexander2000}，\citetitle{alexander2000}，455--458 \\
\ref{sec:AS-2} & 第 I 部分：問題的性質 & \citeauthor{alexander2000}，\citetitle{alexander2000}，459--467 \\
\ref{sec:AS-3} & 第 II 部分：定紛止爭功能及其最適所在 & \citeauthor{alexander2000}，\citetitle{alexander2000}，467--478 \\
\ref{sec:AS-4} & 第 III 部分：最高法院的功能 & \citeauthor{alexander2000}，\citetitle{alexander2000}，478--481 \\
\ref{sec:AS-5} & 結論 & \citeauthor{alexander2000}，\citetitle{alexander2000}，481--482 \\
\addlinespace[6pt]
\multicolumn{3}{@{}l}{Kramer〈Judicial Supremacy and the End of Judicial Restraint〉} \\* \addlinespace[2pt]
\ref{sec:Kramer-1} & 摘要與論題 & \citeauthor{kramer2012}，\citetitle{kramer2012}，621--622 \\
\ref{sec:Kramer-2} & 建國時期的兩種司法審查理論 & \citeauthor{kramer2012}，\citetitle{kramer2012}，622--624 \\
\ref{sec:Kramer-3} & 解釋方法與自制 & \citeauthor{kramer2012}，\citetitle{kramer2012}，624--626 \\
\ref{sec:Kramer-4} & 1800 年至新政 & \citeauthor{kramer2012}，\citetitle{kramer2012}，626--630 \newline \citeauthor{kramer2004}，\citetitle{kramer2004} \\
\ref{sec:Kramer-5} & Brown 案之後的自由派與保守派 & \citeauthor{kramer2012}，\citetitle{kramer2012}，630--632 \\
\ref{sec:Kramer-6} & 解釋理論的興起與自制的衰退 & \citeauthor{kramer2012}，\citetitle{kramer2012}，632--634 \\
\addlinespace[6pt]
\multicolumn{3}{@{}l}{Ferejohn 與 Pasquino〈Constitutional Courts as Deliberative Institutions〉} \\* \addlinespace[2pt]
\ref{sec:FP-1} & 導論：Kelsen 型模式的擴散與合法性負擔 & \citeauthor{ferejohn2002}，\citetitle{ferejohn2002}，21--22 \newline \citeauthor{rawls1993}，\citetitle{rawls1993}，231--236 \\
\ref{sec:FP-2} & 審議的期待 & \citeauthor{ferejohn2002}，\citetitle{ferejohn2002}，22--25 \\
\ref{sec:FP-3} & 民主制度中的審議 & \citeauthor{ferejohn2002}，\citetitle{ferejohn2002}，25--27 \\
\ref{sec:FP-4} & 憲法法院在現代民主中的位置 & \citeauthor{ferejohn2002}，\citetitle{ferejohn2002}，28--33 \\
\ref{sec:FP-5} & 憲法法院的審議 & \citeauthor{ferejohn2002}，\citetitle{ferejohn2002}，33--36 \\
\addlinespace[6pt]
\multicolumn{3}{@{}l}{與前後幾週的關係及課前問題} \\* \addlinespace[2pt]
\ref{sec:relations-1} & 與 W04 至 W08、W10 的對照 & \citeauthor{waldron2006}，\citetitle{waldron2006}，1350 \\
\ref{sec:relations-2} & 課前可以先想的問題 & \citefield{jy185}{title} \newline \citefield{csfa}{title}，第 38 條第 1 項 \newline \citefield{roc-constitution}{title}，第 78 條 \newline \citefield{jy328}{title} \newline \citefield{jy261}{title} \\
\addlinespace[6pt]
\multicolumn{3}{@{}l}{與財政法、稅捐法的關係} \\* \addlinespace[2pt]
\ref{sec:55_fiscal-1} & 一、Alexander 與 Schauer：解釋權威與定紛止爭 & \citeauthor{alexander2000}，\citetitle{alexander2000}，457--463, 476 \newline \citefield{roc-constitution}{title}，第 19 條 \newline \citefield{jy217}{title} \\
\ref{sec:55_fiscal-2} & 二、Kramer：自制的兩種理由與司法至上的興起 & \citeauthor{kramer2012}，\citetitle{kramer2012}，626--630 \\
\ref{sec:55_fiscal-3} & 三、Ferejohn 與 Pasquino：審議型態與立法者角色 & \citeauthor{ferejohn2002}，\citetitle{ferejohn2002}，33 \\
\addlinespace[6pt]
\end{longtable}
\end{small}
